\subsubsection{原函数与常微分方程}\label{subsec:distribution-primitives-odes}

若一个函数的导数为零, 则该函数为常数. 现在为分布建立相应的结果.

\begin{Theorem}
\label{thm:distribution-zero-gradient-constant}
设 $\Omega$ 连通, 且 $u\in\mathcal{D}^{\prime}(\Omega)$ 满足 $\nabla u=0$. 那么 $u$ 为常数.
\end{Theorem}

\begin{Proof}
首先考虑一维情形. 设 $\Omega=I$ 为区间. 条件 $u'=0$ 意味着对每个 $\phi\in\mathcal{D}(I)$ 都有 $(u,\phi')=0$. 换言之, 对每个为某一试验函数之导数的试验函数 $\psi$, 都有 $(u,\psi)=0$. 容易看出, $\psi$ 是试验函数的导数当且仅当 $\int_I\psi(x)\,\mathrm{d}x=0$. 现在设 $\phi_0$ 为任意积分为一的试验函数. 那么任意 $\phi\in\mathcal{D}(I)$ 都可以分解为
\begin{equation}\label{eq:test-function-one-dimensional-zero-integral-decomposition}
\phi(x)=\phi_0(x)\int_I\phi(s)\,\mathrm{d}s+\psi(x),
\end{equation}
其中 $\psi$ 的积分为零. 因此,
\begin{equation}\label{eq:zero-derivative-distribution-constant-pairing}
(u,\phi)=(u,\phi_0)\int_I\phi(x)\,\mathrm{d}x,
\end{equation}
因而 $u$ 等于常数 $(u,\phi_0)$.

接着考虑 $\Omega$ 为区间乘积的情形: $\Omega=(a_1,b_1)\times(a_2,b_2)\times\cdots\times(a_m,b_m)$. 此时, 设 $\phi_i\in\mathcal{D}(a_i,b_i)$ 为积分为一的一维试验函数. 现在将任意 $\phi\in\mathcal{D}(\Omega)$ 分解如下:
\begin{equation}\label{eq:test-function-product-domain-first-decomposition}
\phi(x_1,\ldots,x_m)=\phi_1(x_1)\int_{a_1}^{b_1}\phi(s_1,x_2,\ldots,x_m)\,\mathrm{d}s_1+\psi_1(x_1,\ldots,x_m).
\end{equation}
函数 $\psi_1$ 具有如下性质
\begin{equation}\label{eq:test-function-first-remainder-zero-integral}
\int_{a_1}^{b_1}\psi_1(x_1,x_2,\ldots,x_m)\,\mathrm{d}x_1=0
\end{equation}
对每个 $(x_2,\ldots,x_m)$ 成立; 因而
\begin{equation}\label{eq:test-function-first-remainder-primitive}
\int_{a_1}^{x_1}\psi_1(s,x_2,\ldots,x_m)\,\mathrm{d}s
\end{equation}
仍为试验函数. 由于 $\partial u/\partial x_1=0$, 得到 $(u,\psi_1)=0$. 接着写成
\begin{equation}\label{eq:test-function-product-domain-second-decomposition}
\begin{aligned}
&\phi_1(x_1)\int_{a_1}^{b_1}\phi(s_1,x_2,\ldots,x_m)\,\mathrm{d}s_1\\
&\quad=\phi_1(x_1)\phi_2(x_2)\int_{a_1}^{b_1}\int_{a_2}^{b_2}\phi(s_1,s_2,x_3,\ldots,x_m)\,\mathrm{d}s_2\,\mathrm{d}s_1\\
&\qquad+\phi_1(x_1)\psi_2(x_2,\ldots,x_m),
\end{aligned}
\end{equation}
其中现在有
\begin{equation}\label{eq:test-function-second-remainder-zero-integral}
\int_{a_2}^{b_2}\psi_2(x_2,\ldots,x_m)\,\mathrm{d}x_2=0,
\end{equation}
因而 $(u,\phi_1\psi_2)=0$. 如此继续, 最终得到
\begin{equation}\label{eq:zero-gradient-distribution-product-pairing}
(u,\phi)=(u,\phi_1\phi_2\cdots\phi_m)\int_\Omega\phi(\bm{x})\,\mathrm{d}\bm{x},
\end{equation}
即 $u$ 为常数.

对一般的 $\Omega$, 由刚才证明的结果可知, 每一点都有一个使 $u$ 为常数的邻域, 并且显然, 若两个邻域相交, 则常数必相同 (Problem~\ref{ex:regular-distributions-injective}). 其余部分由 Problem~\ref{ex:distribution-local-vanishing} 得出.
\end{Proof}
